\section{Introduction}
\label{sec:introduction}
Data integration is the process of combining data from multiple heterogeneous sources into a unified, coherent dataset~\cite{doan2012principles,naumann2006dataFusionThreeSteps,khatiwada2022integrating}. This process typically involves multiple steps: schema matching to align attributes across sources, value normalization to standardize data formats, entity matching to identify records referring to the same real-world entity, and data fusion to resolve data conflicts and create a single consolidated record for each entity~\cite{dong2009dataFusion,naumann2006dataFusionThreeSteps}. Traditionally, each step requires substantial manual effort from data engineers who must configure pipeline components and label training data.
Large language models (LLMs) have shown promise in reducing the manual effort required for individual data integration tasks~\cite{freire2025llmDataDiscoveryIntegration}: For schema matching, LLMs can leverage their background knowledge to capture attribute semantics to identify correspondences between heterogeneous schemas~\cite{narayan2022can, zhang2023schema,liu2025magneto}. For entity matching, LLMs have demonstrated competitive performance with BERT-based methods while requiring fewer labeled examples~\cite{peeters2024entity}. These advances raise the question of whether LLMs can successfully generate all artifacts that are required to adapt data integration pipelines to specific use cases, and thus automate the complete data integration process.

This paper investigates to what extent LLMs can replace human involvement in configuring end-to-end data integration pipelines. For this, we present an automatic end-to-end data integration pipeline that relies on GPT-5.2 to perform schema matching, generate value mappings for data normalization, generate training data for entity matching, and generate validation data for selecting conflict resolution heuristics for data fusion. We compare the performance of the automatic pipeline to human-designed pipelines. The human-designed pipelines were created by graduate students as part of a master course on data integration. The students manually labeled training and validation data and configured the pipeline components using the knowledge from the course as well as trial and error. Both the LLM-based and the human-designed pipelines employ methods from the PyDI data integration framework\footnote{\url{https://github.com/wbsg-uni-mannheim/PyDI}}, an open-source library that implements traditional string-based as well as modern embedding- and LLM-based methods for schema matching, value normalization, entity matching, and data fusion. By relying on PyDI, both pipelines can choose scalable methods for tasks that require heavy data processing; the key difference lies in how each pipeline is configured, i.e., which methods are selected and how training and validation data is labeled.

We evaluate the pipelines through three case studies requiring the integration of heterogeneous data describing video games, companies, and music releases. Each use case requires the integration of three data sets having heterogeneous schemas, value formats, and different densities as well as levels of data quality.

% Dataset statistics table for all three use cases
% To be included in Section 2 (Case Studies)

\begin{table*}[t]
\caption{Dataset statistics for the three case studies.}
\label{tab:dataset_statistics}
\centering
\small
\begin{tabular}{llrrcp{9.5cm}}
\toprule
\textbf{Use Case} & \textbf{Source} & \textbf{Rows} & \textbf{Attrs} & \textbf{Density} & \textbf{Attributes} \\
\midrule
\multirow{4}{*}{\textbf{Games}}
 & DBpedia & 46,580 & 6 & 91.0\% & wiki\_ref, title, launch\_yr, studio, system, franchise \\
 & Metacritic & 20,494 & 8 & 97.7\% & mc\_id, game\_title, year\_published, made\_by, console, press\_rating, player\_rating, age\_rating \\
 & Sales & 7,877 & 10 & 98.7\% & rec\_id, prod\_title, launch\_dt, studio, dist, hw, press\_score, comm\_rating, age\_classification, units\_sold\_mm \\
\cmidrule{2-6}
 & \textit{Target Schema} & -- & 12 & -- & id, name, releaseYear, developer, genres, publisher, platform, criticScore, userScore, ESRB, globalSales, series \\
\midrule
\multirow{4}{*}{\textbf{Companies}}
 & Forbes & 2,000 & 7 & 99.2\% & forbes\_url, company, url, region, business\_segment, asset\_value, sales\_figure \\
 & DBpedia & 10,085 & 9 & 66.4\% & entity\_uri, org\_name, established, nation, headquarters, sector, keypeople\_name, total\_assets\_val, annual\_income \\
 & FullContact & 1,931 & 6 & 68.3\% & Attribute\_1, Attribute\_2, Attribute\_3, Attribute\_4, Attribute\_5, Attribute\_6 \\
\cmidrule{2-6}
 & \textit{Target Schema} & -- & 10 & -- & id, name, website, founded, country, city, industry, assets, revenue, founders \\
\midrule
\multirow{4}{*}{\textbf{Music}}
 & Discogs & 22,627 & 8 & 98.4\% & rec\_uid, title\_str, performer, pub\_dt, origin\_loc, imprint, category, tracks\_track-name \\
 & Last.fm & 9,865 & 5 & 78.7\% & item\_code, album\_title, band, tracks\_track-name, album\_length\_min \\
 & MusicBrainz & 4,763 & 7 & 83.3\% & Attribute\_1, Attribute\_2, Attribute\_3, Attribute\_4, Attribute\_5, Attribute\_6, Attribute\_7 \\
\cmidrule{2-6}
 & \textit{Target Schema} & -- & 8 & -- & id, name, artist, release-date, release-country, label, genre, tracks, duration \\
\bottomrule
\end{tabular}
\end{table*}


The contributions of the paper are as follows:

\begin{enumerate}
    \item We present a fully automated end-to-end data integration pipeline that uses an LLM to perform all tasks that traditionally require human expertise, such as choosing methods, labeling training data, and generating validation sets. The pipeline covers the integration steps schema matching, value normalization, entity matching, and data fusion.

    \item We evaluate the automated pipeline along three data integration use cases and compare its results (step-wise and end-to-end) to the performance of human-configured pipelines established by graduate-level data engineers.

    \item We identify limitations of LLM-based methods that we employ and discuss how these limitations could potentially be addressed.

    \item We publish all artifacts of the three case studies (datasets, training, validation and test sets), contributing to the hopefully growing pool of evaluation data for end-to-end data integration. 
\end{enumerate}

The remainder of this paper is organized as follows. Section~\ref{sec:case_studies} describes the three case studies. Sections~\ref{sec:schema_matching} through~\ref{sec:data_fusion} describe how the LLM-based pipeline handles the different integration steps and compares its results to the human-configured pipelines. Section~\ref{sec:end_to_end} presents the end-to-end evaluation. Section~\ref{sec:conclusion} concludes with a discussion of findings and future work. All code and data for reproducing the experiments are available in the accompanying GitHub repository.\footnote{\url{https://github.com/wbsg-uni-mannheim/automatic-data-integration}}
